\documentclass{article}

% Plantilla oficial de NeurIPS (Overleaf: New Project → Templates → "NeurIPS 2025").
% Si la plantilla trae otro año, cambiar el nombre del paquete (p. ej. neurips_2024).
% [preprint] muestra los nombres de los autores sin la nota de "conferencia".
\usepackage[preprint]{neurips_2025}

\usepackage[utf8]{inputenc}
\usepackage[T1]{fontenc}
\usepackage[spanish,es-nodecimaldot,es-noshorthands,es-tabla]{babel}
\usepackage{hyperref}
\usepackage{url}
\usepackage{booktabs}
\usepackage{amsmath,amsfonts}
\usepackage{graphicx}
\usepackage{xcolor}
\usepackage{multirow}

% Si una figura todavía no está en figuras/, se dibuja un recuadro en su lugar (el documento compila igual)
\newcommand{\figura}[2]{%
  \IfFileExists{figuras/#1}{\includegraphics[width=#2]{figuras/#1}}%
  {\fbox{\parbox[c][3cm][c]{#2}{\centering\small Falta \texttt{figuras/#1}}}}}

\title{Mapeo de inundaciones con Sentinel-1 y una U-Net propia:\\
supervisión débil y generalización a Lima 2017}

\author{%
  Integrante 1 \And Integrante 2 \And Integrante 3 \And Integrante 4 \And Integrante 5 \\
  Facultad de Ingeniería Industrial y de Sistemas, Universidad Nacional de Ingeniería, Lima, Perú \\
  Redes Neuronales y Aprendizaje Profundo, 2026-II · Examen parcial, Pregunta 2
}

\begin{document}

\maketitle

\begin{abstract}
Segmentamos píxel a píxel el agua de inundación en imágenes de radar Sentinel-1 (VV, VH) con una U-Net implementada desde cero, entrenada con los chips etiquetados a mano de Sen1Floods11. Un estudio de ablación muestra que la pérdida BCE+Dice supera a BCE y a Dice por separado, y que añadir los 4,160 chips con etiquetas débiles (umbral de Otsu) \emph{empeora} el IoU en validación (0.643 $\rightarrow$ 0.577). Ni las \emph{attention gates} ni entrenar 2.5 veces más mejoran el resultado. El modelo final alcanza un IoU de agua de $0.671 \pm 0.003$ en el test dorado (3 semillas), frente a 0.555 del umbral de Otsu. Fuera de distribución, el rendimiento cae y se vuelve inestable: $0.627 \pm 0.016$ en Bolivia, un evento nunca visto. Construimos además una región de prueba propia para Lima (marzo de 2017) combinando Sentinel-1 del 17 de marzo con la delineación EMSN038 de Copernicus. Allí el IoU es prácticamente nulo ($\leq 0.009$ para todos los métodos). Mostramos que la causa principal es un desajuste entre lo que mide el sensor y lo que marca la referencia: EMSN038 delinea depósitos de un huaico, no agua visible el día de la imagen. A esto se suman agua de río turbulenta y angosta, y sombras de radar en laderas empinadas. El código es reproducible con un solo comando.
\end{abstract}

% ---------------------------------------------------------------------------------------------
\section{Motivación}

Durante una inundación, saber \emph{dónde está el agua} en las horas siguientes es crítico para la respuesta de emergencia, pero las nubes de tormenta inutilizan a los satélites ópticos justo cuando más se los necesita. El radar de apertura sintética (SAR) de Sentinel-1 \citep{torres2012} emite su propia señal de microondas: atraviesa las nubes y opera de noche. El agua en calma refleja la señal lejos del sensor como un espejo y aparece \emph{oscura} (retrodispersión baja, típicamente $< -20$ dB en VH), lo que permite mapearla.

El Perú tiene un precedente concreto. En marzo de 2017, durante El Niño costero, los ríos Rímac y Chillón se desbordaron y los huaicos afectaron Carapongo, Huachipa, Cajamarquilla, Comas y Carabayllo. El Servicio de Gestión de Emergencias de Copernicus activó el mapeo EMSN038 para delinear la zona afectada \citep{emsn038}.

Este trabajo responde tres preguntas:
(i) ¿qué decisiones de entrenamiento importan para una U-Net de segmentación de agua en SAR con pocos datos etiquetados a mano?;
(ii) ¿conviene sumar etiquetas automáticas, abundantes pero ruidosas?; y
(iii) ¿qué tan bien generaliza el modelo a un lugar que nunca vio, en particular a Lima 2017?
Nuestras contribuciones son:
\begin{itemize}
  \item Una U-Net propia, con \emph{attention gates} opcionales, pérdidas y métricas que ignoran los píxeles sin dato, y un protocolo de selección que decide todo en validación y mide la variación entre semillas.
  \item Una ablación de la pérdida, de la supervisión débil (mezcla y preentrenamiento), de la arquitectura y de la duración del entrenamiento.
  \item Una región de prueba de Lima construida a partir de dos fuentes independientes, con un análisis de por qué el IoU es casi nulo y qué fallos del SAR aparecen.
\end{itemize}

% ---------------------------------------------------------------------------------------------
\section{Trabajos relacionados}

\paragraph{Mapeo de inundaciones con SAR por umbral.}
El método operativo clásico separa agua y tierra con un umbral sobre la retrodispersión, a menudo elegido automáticamente con Otsu \citep{otsu1979}. \citet{martinis2009} lo aplican por sub-bloques para manejar histogramas poco bimodales, y \citet{twele2016} construyen una cadena totalmente automática para Sentinel-1 que combina umbral, lógica difusa y el índice topográfico HAND \citep{nobre2011} para descartar zonas elevadas. Estos métodos fallan cuando el agua no es oscura (vegetación inundada, zonas urbanas) o cuando algo seco sí lo es (sombras, superficies lisas).

\paragraph{Segmentación con redes profundas.}
La U-Net \citep{ronneberger2015} combina un codificador que resume el contexto con un decodificador que recupera el detalle mediante conexiones de salto, y es la arquitectura de referencia en segmentación con pocos datos. La Attention U-Net \citep{oktay2018} filtra cada conexión de salto con una compuerta aprendida, y UNet++ \citep{zhou2018} rediseña los saltos con caminos anidados. La pérdida Dice \citep{milletari2016} optimiza directamente el solapamiento y es robusta al desbalance de clases.

\paragraph{Conjuntos de datos y supervisión débil.}
Sen1Floods11 \citep{bonafilia2020} ofrece 446 chips Sentinel-1 etiquetados a mano y 4,385 con etiquetas automáticas, sobre 11 eventos, y reporta que entrenar con las etiquetas débiles puede competir con las manuales. WorldFloods \citep{mateogarcia2021} entrena redes de segmentación de agua a escala global con imágenes ópticas. El aprendizaje con etiquetas ruidosas es un área amplia \citep{song2022}: las redes pueden memorizar el ruido, pero también filtrarlo parcialmente cuando el ruido es sistemático. Este trabajo mide cuál de los dos efectos domina en Sen1Floods11.

% ---------------------------------------------------------------------------------------------
\section{Datos}

\subsection{Sen1Floods11}
Usamos los chips etiquetados a mano (``dorados'') de Sen1Floods11 v1.1: imágenes Sentinel-1 GRD de $512 \times 512$ píxeles a 10 m (bandas VV y VH en dB, EPSG:4326) con máscaras de 1 (agua), 0 (no agua) y $-1$ (sin dato). Una exploración previa a modelar llevó a tres decisiones:
\begin{itemize}
  \item \textbf{Exclusiones.} Excluimos 7 chips inservibles: 6 de Ghana con $\geq 99\%$ de píxeles $-1$, y \texttt{Paraguay\_34417}, con la imagen 100\% vacía. Las particiones oficiales quedan en 250 / 86 / 88 / 15 (train / valid / test / Bolivia).
  \item \textbf{Píxeles sin radar.} Los píxeles sin medición de radar (NaN) pasan a $-1$ en la máscara, porque no se pueden predecir. Además, el $-1$ no está declarado como \emph{nodata} en los GeoTIFF, así que lo tratamos explícitamente al leer.
  \item \textbf{Desbalance.} Solo $\approx 10.6\%$ de los píxeles es agua (mediana de 2.5\% por chip, y 52 chips sin agua). Por eso usamos métricas de la clase agua y no la exactitud global.
\end{itemize}
La partición Bolivia contiene un evento ausente de train, y la usamos como prueba de generalización geográfica dentro del propio dataset.

\textbf{Etiquetas débiles.} Los chips débiles (\texttt{S1Weak} con \texttt{S1OtsuLabelWeak}) se generan con un umbral de Otsu ajustado por evento. Para evitar fuga de datos aplicamos dos filtros:
\begin{itemize}
  \item Excluimos los 224 chips débiles de Bolivia.
  \item Excluimos todo chip débil que se solape en más de 1\% de su área con un chip de valid, test o Bolivia.
\end{itemize}
Quedan 4,160 chips de 11 eventos, con un porcentaje medio de agua por chip de entre 4.6\% (EE.~UU.) y 30.5\% (Pakistán).

\subsection{Región de prueba de Lima (construida por el grupo)}

\textbf{Imagen SAR.} Usamos la escena Sentinel-1B del 17 de marzo de 2017 (órbita descendente, track 171, VV+VH). Es la más cercana al pico del evento (16--19 de marzo) con ambas polarizaciones: las pasadas del 18 y del 24 de marzo solo tienen VV. La procesamos con Google Earth Engine \citep{gorelick2017} desde la colección \texttt{COPERNICUS/S1\_GRD}, con calibración, corrección de terreno y paso a dB. Es la misma cadena con la que se construyó Sen1Floods11. La exportamos a 10 m en EPSG:4326 sobre el rectángulo $[-77.20, -12.10, -76.80, -11.75]$, que cubre los ríos Rímac y Chillón: son $4454 \times 3898$ píxeles, de los cuales 17.8\% no tiene dato (el borde de la franja, casi todo mar).

\textbf{Etiquetas EMSN038.} Leímos el \emph{geodatabase} del producto (EPSG:32618), lo reproyectamos a EPSG:4326 y lo rasterizamos sobre la misma grilla de la imagen. Marcamos con $-1$ todo píxel fuera del área analizada por Copernicus (\texttt{EMSN038\_AOI}) o sin dato de radar. La capa de inundación \texttt{Potential\_Maximum\_flood} contiene dos tipos de polígono, que evaluamos por separado: el \emph{río} (2 polígonos) y la extensión máxima del \emph{huaico} (11 polígonos). Como referencia adicional usamos el cauce de los ríos de la cartografía base (\texttt{WatrcrsA}). Dentro del área evaluable (13.5 millones de píxeles), la huella completa ocupa solo 1.02\% y el cauce, 0.36\%.

\textbf{Advertencia metodológica.} Estas etiquetas no se generaron con el criterio de Sen1Floods11. Se delinearon sobre imágenes ópticas WorldView-2 del 18 y 21 de marzo, posteriores a la imagen SAR, y su reporte técnico indica que \emph{no se detectó agua estancada con Sentinel-1}: el agua se retiró en horas y lo que se delinea son huellas del evento (lodo, humedad, daños). Son, por lo tanto, una supervisión débil propia y no una verdad de terreno.

% ---------------------------------------------------------------------------------------------
\section{Metodología}

\textbf{Preprocesamiento.} Cada banda se recorta a sus percentiles 0.1 y 99.9 y se estandariza con la media y la desviación estimadas \emph{solo con train} (20,000 píxeles por chip). Los NaN pasan a 0 (y ya son $-1$ en la máscara). Las mismas estadísticas se aplican sin cambios a test, Bolivia y Lima. Al entrenar aplicamos espejos, rotaciones de $90^\circ$ y un recorte aleatorio de $256 \times 256$. No modificamos el brillo, porque en el radar el brillo es la señal. La evaluación usa el chip completo de 512.

\textbf{Modelo.} Implementamos una U-Net \citep{ronneberger2015} de 4 niveles con filtros $(32, 64, 128, 256)$ y un cuello de 512. Cada bloque tiene dos convoluciones $3\times3$ con BatchNorm y ReLU, la bajada es por \emph{max-pooling} y la subida por convolución transpuesta. La entrada es de 2 canales y la salida, de 1 \emph{logit}; el modelo suma $\approx 7.8$ M de parámetros (la mitad de filtros que la U-Net original, porque hay solo 250 chips). La variante Attention U-Net \citep{oktay2018} filtra cada salto $x$ con el contexto del decodificador $g$:
\begin{equation}
  \alpha = \sigma\!\big(\psi^\top \operatorname{ReLU}(W_x x + W_g g)\big), \qquad \hat{x} = \alpha \odot x,
\end{equation}
con una capa intermedia de la mitad de canales del salto, lo que añade un 1.1\% de parámetros.

\textbf{Pérdidas con máscara.} Sea $v_i = \mathbb{1}[y_i \neq -1]$ y $p_i = \sigma(z_i)$. Entonces:
\begin{equation}
  \mathcal{L}_\text{BCE} = \frac{\sum_i v_i\,\ell_\text{BCE}(z_i, y_i)}{\sum_i v_i}, \qquad
  \mathcal{L}_\text{Dice} = 1 - \frac{2 \sum_i v_i p_i y_i + 1}{\sum_i v_i p_i + \sum_i v_i y_i + 1}, \qquad
  \mathcal{L}_\text{BCE+Dice} = \tfrac12 \mathcal{L}_\text{BCE} + \tfrac12 \mathcal{L}_\text{Dice}.
\end{equation}
Una prueba unitaria verifica que cambiar la predicción en un píxel $-1$ no altera ninguna de las tres pérdidas.

\textbf{Métricas.} Acumulamos la matriz de confusión de la clase agua sobre todos los píxeles válidos de una partición y reportamos IoU $= TP/(TP+FP+FN)$ (métrica principal), F1, precisión, recall, omisión ($FN/(TP+FN)$) y comisión ($FP/(TP+FP)$). También reportamos el IoU promedio por chip, la métrica de los autores de Sen1Floods11, que pesa igual a cada chip.

\textbf{Entrenamiento y selección.} Entrenamos con AdamW (lr $10^{-3}$, \emph{weight decay} $10^{-4}$), \emph{learning rate} coseno, 40 épocas, lotes de 8 y precisión mixta, en una GPU T4 ($\approx 5$ s por época). Conservamos el mejor \emph{checkpoint} por IoU en valid. El protocolo de selección es estricto:
\begin{itemize}
  \item \textbf{Solo valid decide.} Valid elige la pérdida, la arquitectura, la duración y el \emph{checkpoint}.
  \item \textbf{Regla de mejora.} Un cambio se adopta solo si mejora el IoU de valid en más de 0.02, porque la variación entre corridas es $\approx 0.01$.
  \item \textbf{Modelo final.} Con todo decidido, la receta ganadora se reentrena con train+valid (336 chips) durante el mismo número fijo de épocas.
  \item \textbf{Una sola mirada.} Test, Bolivia y Lima se evalúan una sola vez, al final.
  \item \textbf{Semillas.} Ambas recetas (``solo train'' y ``train+valid'') se repiten con 3 semillas (42, 1, 2), y reportamos media $\pm$ desviación.
\end{itemize}

\textbf{Línea base.} En Sen1Floods11 usamos las máscaras \texttt{S1OtsuLabelHand} del propio dataset (umbral de Otsu ajustado por evento). En Lima aplicamos Otsu sobre el histograma de VH de la escena.

% ---------------------------------------------------------------------------------------------
\section{Experimentos y resultados}

\subsection{Estudio de ablación (validación)}

\begin{table}[t]
  \caption{Ablación en valid (86 chips). E1--E5 varían la pérdida y la supervisión; E6 y E7 parten de E3 y cambian una sola cosa. ``Últ.~5'': IoU medio de las últimas 5 épocas (estabilidad).}
  \label{tab:ablacion}
  \centering
  \small
  \begin{tabular}{llcccc}
    \toprule
    Exp. & Decisión & IoU & Precisión & Recall & Últ.~5 \\
    \midrule
    Otsu & umbral clásico (línea base) & 0.507 & 0.634 & 0.716 & -- \\
    E1 & U-Net, BCE, solo dorados & 0.623 & 0.851 & 0.700 & 0.605 \\
    E2 & U-Net, Dice, solo dorados & 0.640 & 0.807 & 0.756 & 0.633 \\
    E3 & U-Net, BCE+Dice, solo dorados & \textbf{0.643} & 0.787 & \textbf{0.778} & 0.618 \\
    \midrule
    E4 & E3 + 500 débiles por época (mezcla) & 0.577 & 0.779 & 0.690 & 0.567 \\
    E5a & solo débiles (20 épocas) & 0.548 & 0.768 & 0.657 & 0.542 \\
    E5 & E5a $\rightarrow$ ajuste con dorados (lr $3\cdot10^{-4}$) & 0.620 & 0.792 & 0.741 & 0.601 \\
    \midrule
    E6 & E3 con Attention U-Net & 0.643 & 0.810 & 0.757 & 0.618 \\
    E7 & E3 con 100 épocas & 0.639 & 0.835 & 0.731 & 0.626 \\
    \bottomrule
  \end{tabular}
\end{table}

La Tabla~\ref{tab:ablacion} resume la ablación.
\begin{itemize}
  \item \textbf{Pérdida (b).} Las pérdidas con Dice suben el recall de 0.70 a 0.76--0.78 a cambio de algo de precisión, y son más estables: con BCE, el IoU de valid cae a 0.13 en la época 17. E2 y E3 empatan (diferencia de 0.003), y elegimos BCE+Dice por su mayor recall. En una emergencia, no ver agua real es más grave que marcar un poco de más.
  \item \textbf{Etiquetas débiles (a).} No ayudan: mezclarlas resta 0.066 de IoU (E4) y preentrenar con ellas resta 0.022 (E5). Al inspeccionar las máscaras débiles vemos que son manchas suavizadas, sin ríos finos ni charcos pequeños. E4 aprende ese estilo y pierde ríos estrechos, y además hereda el error de Otsu sobre campos lisos. Aun así, E5a, entrenado \emph{solo} con etiquetas de Otsu, supera al propio Otsu (0.548 frente a 0.507, con comisión de 0.23 frente a 0.37): la red no copia el umbral, sino que aprende forma y contexto, y filtra parte del ruido sistemático \citep{song2022}.
  \item \textbf{Arquitectura y duración.} Las \emph{attention gates} (E6) empatan con E3 y solo desplazan el balance hacia la precisión. Entrenar 100 épocas (E7) estabiliza las curvas, pero no mejora el mejor IoU. Concluimos que con 250 chips el límite son los datos y no la capacidad del modelo, así que el modelo elegido es \textbf{E3}.
\end{itemize}

\subsection{Evaluación en distribución (test)}

\begin{table}[t]
  \caption{Test (88 chips de los eventos de entrenamiento) y Bolivia (15 chips, evento no visto). Media $\pm$ desviación de 3 semillas, salvo Otsu. F1 por píxel $= 2\,\mathrm{IoU}/(1+\mathrm{IoU})$.}
  \label{tab:test}
  \centering
  \small
  \begin{tabular}{lcccc}
    \toprule
    & IoU test & Recall test & IoU Bolivia & Recall Bolivia \\
    \midrule
    Otsu (umbral por evento) & 0.555 & 0.748 & \textbf{0.718} & \textbf{0.870} \\
    U-Net E3, solo train & $0.665 \pm 0.002$ & $0.793 \pm 0.021$ & $0.669 \pm 0.029$ & $0.744 \pm 0.047$ \\
    U-Net E3, train+valid (final) & $\mathbf{0.671 \pm 0.003}$ & $0.780 \pm 0.004$ & $0.627 \pm 0.016$ & $0.680 \pm 0.021$ \\
    \bottomrule
  \end{tabular}
\end{table}

En test, el modelo final (semilla 42) obtiene IoU 0.670, F1 0.802, precisión 0.831 y recall 0.776, frente a IoU 0.555 y F1 0.714 de Otsu (Tabla~\ref{tab:test}). La mejora de +0.12 de IoU es muy estable entre semillas (desviación de 0.003). La matriz de confusión en píxeles es la siguiente:
\begin{center}
\small
\begin{tabular}{lrr}
  \toprule
  & Predice no agua & Predice agua \\
  \midrule
  Real: no agua & 17,532,363 & 396,588 \\
  Real: agua & 563,184 & 1,951,824 \\
  \bottomrule
\end{tabular}
\end{center}
La red deja sin ver más agua de la que inventa: 563 mil píxeles de omisión frente a 397 mil de comisión. Es un modelo algo conservador. Por evento, la red supera a Otsu en los 10 eventos de test. Los mejores resultados aparecen donde hay agua abierta y extensa (Nigeria 0.89, Sri Lanka 0.88, Mekong 0.86) y los peores donde el agua es difusa, fina o se confunde con suelo liso (Pakistán 0.19, Somalia 0.42, Ghana 0.53). El porcentaje de píxeles sin dato no explica el ranking: Pakistán y Nigeria tienen ambos $\approx 28$--$30\%$. El IoU promedio por chip, la métrica de los autores, es de 0.400 en test. Es mucho menor que el IoU por píxel porque los chips con poca agua obtienen IoU bajos y pesan igual que los demás.

\subsection{Generalización geográfica: Bolivia y Lima}

\textbf{Bolivia.} El IoU baja de 0.671 a $0.627 \pm 0.016$ y la desviación entre semillas es entre 5 y 15 veces la de test. Además, sumar valid al entrenamiento \emph{mejora} el test (+0.007, consistente en las 3 semillas) pero \emph{empeora} Bolivia ($-0.042$, también consistente). Nuestra hipótesis es que valid comparte eventos con train, así que entrenar con él especializa a la red en esos paisajes. Otsu gana en Bolivia (0.718), pero su umbral se ajustó por evento (\texttt{VH\_thresh} entre $-19.6$ y $-25.1$ dB), mientras que la red nunca vio esa escena. La matriz de confusión del modelo final en Bolivia muestra 4 veces más omisión que comisión (151 mil frente a 37 mil píxeles).

\begin{table}[t]
  \caption{Lima, 17 de marzo de 2017, dentro del área analizada por EMSN038. Media de 3 semillas (desviación $\leq 0.001$ en IoU). ``\% agua'' es el porcentaje del área que cada método marca como agua.}
  \label{tab:lima}
  \centering
  \small
  \begin{tabular}{llccccc}
    \toprule
    Referencia & Método & IoU & Recall & Comisión & IoU/chip & \% agua \\
    \midrule
    \multirow{3}{*}{Huella completa (1.02\%)} & Otsu & 0.009 & 0.425 & 0.991 & -- & 49.5 \\
     & U-Net, solo train & 0.003 & 0.026 & 0.997 & 0.005 & 9.2 \\
     & U-Net, train+valid & 0.002 & 0.010 & 0.998 & 0.002 & 4.2 \\
    \midrule
    \multirow{3}{*}{Río (EMSN038)} & Otsu & 0.006 & 0.676 & 0.994 & -- & 49.5 \\
     & U-Net, solo train & 0.003 & 0.053 & 0.997 & -- & 9.2 \\
     & U-Net, train+valid & 0.002 & 0.022 & 0.998 & -- & 4.2 \\
    \midrule
    \multirow{3}{*}{Huaico (EMSN038)} & Otsu & 0.002 & 0.212 & 0.998 & -- & 49.5 \\
     & U-Net, solo train & 0.000 & 0.003 & 1.000 & -- & 9.2 \\
     & U-Net, train+valid & 0.000 & 0.000 & 1.000 & -- & 4.2 \\
    \midrule
    \multirow{3}{*}{Cauce (cartografía base)} & Otsu & 0.005 & 0.708 & 0.995 & -- & 49.5 \\
     & U-Net, solo train & 0.002 & 0.053 & 0.998 & 0.005 & 9.2 \\
     & U-Net, train+valid & 0.001 & 0.019 & 0.999 & 0.002 & 4.2 \\
    \bottomrule
  \end{tabular}
\end{table}

\textbf{Lima.} El IoU es prácticamente nulo para todos los métodos y todas las referencias (Tabla~\ref{tab:lima}). La brecha frente al test pasa de 0.67 a $\leq 0.003$, y el IoU por chip, de 0.400 a $\leq 0.005$. Antes de culpar a la red verificamos el \emph{pipeline}. Los polígonos de EMSN038 se alinean con los ríos visibles en VV, y el mar, que está fuera del área evaluada, se detecta completo como agua (Figura~\ref{fig:lima}). La normalización y la red funcionan, por lo que el problema está en qué miden la referencia y el sensor. Proponemos cuatro hipótesis sobre la causa:

\begin{enumerate}
  \item \textbf{Desajuste temporal y de concepto entre etiqueta e imagen.} La huella del huaico (11 polígonos, en su mayoría en Carapongo y Huachipa) cae sobre zonas grises de VV, sin la firma oscura del agua. Es lodo y suelo removido, delineado en imágenes ópticas de 1 a 4 días después de la imagen SAR. Las 6 redes obtienen recall de 0.000 a 0.003 contra el huaico. Es el resultado esperado, y coincide con el reporte de Copernicus, que tampoco detectó agua estancada con Sentinel-1.
  \item \textbf{Agua de río turbulenta y angosta.} Incluso contra el río la red ve solo el 2--5\% del agua. Durante El Niño costero el Rímac y el Chillón llevaban caudales altos y cargados de sedimento. Esa superficie rugosa devuelve eco y no se ve oscura, y además el cauce mide pocos píxeles de ancho, que es el mismo fallo de ``ríos finos'' observado en test. Los dos polígonos de río de EMSN038 y el cauce de la cartografía base se comportan igual, así que el resultado no depende de quién dibujó el río.
  \item \textbf{Relieve y geometría de adquisición.} Casi toda el agua que marca la red está en laderas de los cerros que rodean los valles (Figura~\ref{fig:lima}). Con la geometría de órbita descendente y el ángulo de incidencia local que impone la pendiente, las caras opuestas al sensor quedan en \emph{sombra de radar} y se ven tan oscuras como el agua. Sen1Floods11 casi no contiene relieve andino, así que la red nunca aprendió a distinguirlas. Las 6 redes coinciden solo en el 20.7\% de los píxeles que alguna marca como agua: fuera de distribución, la respuesta depende del azar del entrenamiento, igual que en Bolivia.
  \item \textbf{Cobertura urbana.} La ciudad aparece muy brillante por el doble rebote entre suelo y edificios, un paisaje ausente en Sen1Floods11. La red casi no marca agua en la ciudad: el agua urbana, si la hubo, no se vería oscura.
\end{enumerate}

Otsu no es una alternativa en Lima. El histograma de VH es unimodal, porque casi no hay agua en la escena, así que el umbral queda en $-17.7$ dB, fuera del rango de Sen1Floods11, y marca como agua la mitad del área. Su recall aparentemente alto frente al río (0.68) se debe solo a eso.

\begin{figure}[t]
  \centering
  \figura{lima_escena.png}{0.48\linewidth}\hfill
  \figura{lima_errores_1.png}{0.48\linewidth}
  \caption{Lima, 17 de marzo de 2017. Izquierda: agua según la red final (rojo) y huella de EMSN038 (cian). El mar se detecta, pero la mayor parte de la comisión está en laderas en sombra de radar. Derecha: Rímac (Huachipa--Carapongo). Naranja indica agua de la referencia no detectada, rojo agua inventada y azul aciertos.}
  \label{fig:lima}
\end{figure}

\subsection{Análisis cualitativo de fallos del SAR}

Calculamos el IoU de cada chip de test y Bolivia, y analizamos los mapas de error de los peores (Figura~\ref{fig:fallos}) junto con la escena de Lima. Identificamos cuatro tipos de error característicos del SAR:
\begin{enumerate}
  \item \textbf{Agua que no se ve oscura} (omisión). En la vegetación inundada, los tallos sobre el agua producen doble rebote y la superficie se ve gris o brillante. Lo mismo ocurre con el agua somera, rugosa o turbulenta, como el Rímac en 2017. Ejemplos: \texttt{Ghana\_313799}, \texttt{Somalia\_699062}, \texttt{Bolivia\_312675}.
  \item \textbf{Superficies lisas que no son agua} (comisión). Los campos de cultivo, el suelo húmedo o desnudo y las pistas reflejan como un espejo. Ejemplos: \texttt{Pakistan\_70625}, donde la comisión del evento llega a 0.55, y \texttt{Paraguay\_76868}.
  \item \textbf{Ríos y canales finos} (omisión). Con 1--3 píxeles de ancho a 10 m, el río se mezcla con la orilla. Ejemplos: \texttt{USA\_670826} y los ríos de Lima.
  \item \textbf{Sombra de radar} (comisión). Las laderas de espaldas al sensor no reciben señal. Es el error dominante en Lima (Figura~\ref{fig:lima}).
\end{enumerate}

\begin{figure}[t]
  \centering
  \figura{peores_chips.png}{0.9\linewidth}
  \caption{Chips de test y Bolivia con menor IoU (al menos 1\% de agua): VV, VH y mapa de errores del modelo final. Azul indica acierto, rojo agua inventada (comisión), naranja agua no vista (omisión) y gris sin dato.}
  \label{fig:fallos}
\end{figure}

% ---------------------------------------------------------------------------------------------
\section{Discusión y limitaciones}

\begin{itemize}
  \item \textbf{Las etiquetas de Lima no son verdad de terreno.} EMSN038 delinea la huella del evento en imágenes ópticas posteriores, no el agua visible en la fecha del SAR. Un IoU cercano a cero mide sobre todo ese desajuste, y no permite estimar cuánto agua real se pierde. La conclusión válida es cualitativa: el modelo, y cualquier método basado en ``agua oscura'', no sirve para huaicos de agua rápida como los de Lima 2017.
  \item \textbf{La generalización es frágil y ruidosa.} Con solo 15 chips, Bolivia muestra desviaciones de hasta 0.03 entre semillas. Una sola corrida habría llevado a conclusiones distintas sobre si sumar valid ayuda. Reportar semillas es imprescindible fuera de distribución.
  \item \textbf{La comparación con Otsu en Bolivia está sesgada a favor de Otsu}, porque su umbral se ajustó por evento.
  \item \textbf{El resultado sobre las etiquetas débiles depende de su calidad.} Las máscaras de Otsu de Sen1Floods11 son suavizadas y comparten los errores del umbral. Con otras fuentes débiles (por ejemplo, Sentinel-2) o con pérdidas robustas al ruido, el resultado podría cambiar.
  \item \textbf{Trabajo futuro.} Tres extensiones directas son: (i) una máscara física de terreno con pendiente o HAND \citep{nobre2011,twele2016} que descarte sombras de radar, decidida a priori y no ajustada a Lima; (ii) detección de cambios entre una imagen previa y una posterior, más adecuada para huaicos; y (iii) la fusión con bandas de Sentinel-2 (el bono del enunciado), que no se exploró.
  \item \textbf{Reproducibilidad.} La precisión mixta y cuDNN no son totalmente deterministas en GPU, así que los números pueden variar $\approx 0.01$ entre ejecuciones.
\end{itemize}

\section{Conclusiones}

Una U-Net propia de 7.8 M de parámetros, entrenada solo con 250 chips etiquetados a mano y con la pérdida BCE+Dice, supera al umbral de Otsu por +0.12 de IoU dentro de distribución, de forma estable entre semillas. Las etiquetas débiles de Otsu, las \emph{attention gates} y un entrenamiento más largo no aportan: el cuello de botella son los datos etiquetados. Fuera de distribución el rendimiento cae y se vuelve variable. En Lima 2017 el modelo no recupera la delineación de EMSN038, porque el evento fue un huaico cuyo rastro no tiene la firma de agua oscura que el SAR (y la red) necesitan. A eso se suman ríos turbulentos y angostos, y sombras de radar en un relieve ausente del entrenamiento. Para un despliegue en el Perú, un modelo de este tipo necesitaría ejemplos locales con relieve andino, una máscara de terreno y etiquetas tomadas en la misma fecha que la imagen SAR.

\paragraph{Código.} Repositorio: \url{https://github.com/USUARIO/REPO}. Todo el estudio se reproduce con \texttt{bash run\_all.sh} (o \texttt{SIN\_DEBILES=1 bash run\_all.sh}, sin etiquetas débiles).

% ---------------------------------------------------------------------------------------------
\small
\begin{thebibliography}{99}

\bibitem[Bonafilia et~al.(2020)]{bonafilia2020}
D.~Bonafilia, B.~Tellman, T.~Anderson y E.~Issenberg.
\newblock Sen1Floods11: A georeferenced dataset to train and test deep learning flood algorithms for Sentinel-1.
\newblock En \emph{CVPR Workshops}, 2020.

\bibitem[Copernicus EMS(2017)]{emsn038}
Copernicus Emergency Management Service.
\newblock EMSN038: Floods in Lima, Peru. Reporte técnico y geodatabase.
\newblock \url{https://mapping.emergency.copernicus.eu/activations/EMSN038/}, 2017.

\bibitem[Gorelick et~al.(2017)]{gorelick2017}
N.~Gorelick, M.~Hancher, M.~Dixon, S.~Ilyushchenko, D.~Thau y R.~Moore.
\newblock Google Earth Engine: Planetary-scale geospatial analysis for everyone.
\newblock \emph{Remote Sensing of Environment}, 202:18--27, 2017.

\bibitem[Martinis et~al.(2009)]{martinis2009}
S.~Martinis, A.~Twele y S.~Voigt.
\newblock Towards operational near real-time flood detection using a split-based automatic thresholding procedure on high resolution TerraSAR-X data.
\newblock \emph{Natural Hazards and Earth System Sciences}, 9:303--314, 2009.

\bibitem[Mateo-Garcia et~al.(2021)]{mateogarcia2021}
G.~Mateo-Garcia, J.~Veitch-Michaelis, L.~Smith, S.~V. Oprea, G.~Schumann, Y.~Gal, A.~G. Baydin y D.~Backes.
\newblock Towards global flood mapping onboard low cost satellites with machine learning.
\newblock \emph{Scientific Reports}, 11:7249, 2021.

\bibitem[Milletari et~al.(2016)]{milletari2016}
F.~Milletari, N.~Navab y S.-A. Ahmadi.
\newblock V-Net: Fully convolutional neural networks for volumetric medical image segmentation.
\newblock En \emph{3DV}, 2016.

\bibitem[Nobre et~al.(2011)]{nobre2011}
A.~D. Nobre, L.~A. Cuartas, M.~Hodnett, C.~D. Rennó, G.~Rodrigues, A.~Silveira, M.~Waterloo y S.~Saleska.
\newblock Height Above the Nearest Drainage -- a hydrologically relevant new terrain model.
\newblock \emph{Journal of Hydrology}, 404:13--29, 2011.

\bibitem[Oktay et~al.(2018)]{oktay2018}
O.~Oktay, J.~Schlemper, L.~Le~Folgoc, M.~Lee, M.~Heinrich, K.~Misawa, K.~Mori, S.~McDonagh, N.~Y. Hammerla, B.~Kainz, B.~Glocker y D.~Rueckert.
\newblock Attention U-Net: Learning where to look for the pancreas.
\newblock En \emph{MIDL}, 2018.

\bibitem[Otsu(1979)]{otsu1979}
N.~Otsu.
\newblock A threshold selection method from gray-level histograms.
\newblock \emph{IEEE Transactions on Systems, Man, and Cybernetics}, 9(1):62--66, 1979.

\bibitem[Ronneberger et~al.(2015)]{ronneberger2015}
O.~Ronneberger, P.~Fischer y T.~Brox.
\newblock U-Net: Convolutional networks for biomedical image segmentation.
\newblock En \emph{MICCAI}, 2015.

\bibitem[Song et~al.(2022)]{song2022}
H.~Song, M.~Kim, D.~Park, Y.~Shin y J.-G. Lee.
\newblock Learning from noisy labels with deep neural networks: A survey.
\newblock \emph{IEEE Transactions on Neural Networks and Learning Systems}, 2022.

\bibitem[Torres et~al.(2012)]{torres2012}
R.~Torres et~al.
\newblock GMES Sentinel-1 mission.
\newblock \emph{Remote Sensing of Environment}, 120:9--24, 2012.

\bibitem[Twele et~al.(2016)]{twele2016}
A.~Twele, W.~Cao, S.~Plank y S.~Martinis.
\newblock Sentinel-1-based flood mapping: a fully automated processing chain.
\newblock \emph{International Journal of Remote Sensing}, 37(13):2990--3004, 2016.

\bibitem[Zhou et~al.(2018)]{zhou2018}
Z.~Zhou, M.~M. Rahman~Siddiquee, N.~Tajbakhsh y J.~Liang.
\newblock UNet++: A nested U-Net architecture for medical image segmentation.
\newblock En \emph{DLMIA (MICCAI Workshops)}, 2018.

\end{thebibliography}

\end{document}
